\begin{latin}
\begin{center}\textbf{Abstract}\end{center}
Adversarial training improves the robustness of neural networks but incurs substantial computational cost when adversarial examples are generated using iterative attacks. Fast single-step alternatives can suffer from catastrophic overfitting: accuracy against multi-step attacks abruptly collapses despite high single-step accuracy. This thesis studies inexpensive prediction and mitigation of this failure through the geometry of the input loss landscape.

First, epsilon overfitting characterizes the model's sensitivity to restricted perturbation magnitudes and directions. Second, perturbation alignment (PertAlign) measures the cosine similarity between input gradients before and after the attack. A local Taylor analysis relates its decrease to the component of the Hessian-transformed attack direction orthogonal to the initial gradient. Reusing gradients from adversarial-example generation and parameter-update backpropagation enables inexpensive monitoring without explicitly computing a Hessian.

These insights motivate Second-Order Adaptive (SORA) training. SORA uses a directional curvature estimate, exponential averaging, and elementwise random step sizes to adapt its training perturbations. Published experiments cover natural and biomedical image datasets and multiple architectures with fixed method-specific hyperparameters. For PreActResNet-18 under AutoAttack at an evaluation budget of $8/255$, the reported robust accuracies are $32.91\%$ on PathMNIST and $18.56\%$ on ImageNet-100. The thesis develops and contextualizes results from the joint ICML 2026 paper rather than reporting newly executed training experiments. It also discusses the local nature of the approximation, the distinct training and evaluation perturbation sets, and the need for statistical calibration of early warnings.

\bigskip\noindent\textbf{Keywords:}
Fast adversarial training, catastrophic overfitting, epsilon overfitting, perturbation alignment, second-order optimization, adversarial robustness.
\end{latin}
